Con el crecimiento del mercado de los servicios de infraestructura como código~\cite{grandviewresearchInfrastructureCodeMarket}, se ha aprendido que el poder definir el estado de sistemas completos a través de código trae ventajas importantes en velocidad de desarrollo, escalabilidad y costos~\cite{pandyaIntroductionInfrastructureCode2022}. La misma idea también se ha aplicado a entornos de desarrollo~\cite{ghanbariUsingDevelopmentEnvironment2026} o flujos de despliegue continuo~\cite{wesselGitHubActionsImpact2023}; todo lo anterior sugiere que herramientas que permiten definir entidades o procesos de forma declarativa pueden facilitar el ciclo de desarrollo de \emph{software}. En el manejo de paquetes, el panorama es fragmentado: habiendo alternativas por lenguaje, o indirectas como Docker~\cite{Zwinger2026}; y una solución declarativa de propósito general no es del todo clara. Nix lleva años intentando llenar ese espacio --- y las cifras sugieren que está haciendo algo bien: al tener uno de los repositorios de paquetes más grandes de Linux, con 115 mil paquetes~\cite{marakasovRepositoryStatistics}, y un crecimiento del 264\% en número de mantenedores en los últimos seis años~\cite{gg-solutionsLinuxSilentTech2026}.

Por medio de su lenguaje de configuración, Nix permite describir la construcción, instalación y composición de paquetes de \emph{software}~\cite{Dolstra2006}, habilidad que se ha mostrado aplicable en configuración de ambientes de Computación de Alto Rendimiento~\cite{guilloteauPainlessTranspositionReproducible2022,Gomez2020}, sistemas operativos~\cite{Thiberg2025}, despliegue de \emph{software}~\cite{VanDerBurg2014}, orquestación de servicios~\cite{FloxKubernetesUncontained} o entornos de desarrollo~\cite{replitReplitHowWe2021}.

No obstante, a pesar de su versatilidad, Nix presenta una barrera de entrada considerable. Reportes sugieren una curva de aprendizaje pronunciada~\cite{NixCommunitySurvey2024}, atribuida en parte a la complejidad de su lenguaje de configuración (Nixlang) para ciertos usuarios~\cite{fricklerhandwerk2022}, aspecto que también su creador ha identificado como susceptible de mejora~\cite{Dolstra2018}. Estas dificultades podrían estar incidiendo en su adopción relativamente limitada frente a herramientas como Docker~\cite{stackoverflowMostPopularTechnologies}.

Nixlang es la principal interfaz para interactuar con el ecosistema Nix, es un lenguaje de dominio específico (DSL por sus siglas en inglés) diseñado directamente para expresar los constructos de la herramienta, y es, en gran medida, responsable de la flexibilidad que la caracteriza~\cite{NixdevDocumentation}. Sin embargo, esta misma especialización introduce complejidades que afectan su accesibilidad y usabilidad~\cite{Dolstra2018}. Como respuesta, la comunidad ha explorado diversas estrategias para mitigar estas limitaciones, como extensiones al lenguaje, con la incorporación de tipado estático ---esfuerzos que han sido abandonados debido a su complejidad técnica---~\cite{caddetNixNickel,hufschmittCurrentStatePtyx}; agentes de inteligencia artificial para generar configuraciones, aún sin validación empírica sólida en términos de usabilidad~\cite{Schwaighofer2026}; y, en un enfoque más radical, la sustitución del lenguaje por Guile, un lenguaje de propósito general~\cite{Courts2013}, aunque tampoco es un lenguaje muy conocido~\cite{stackoverflowMostPopularTechnologies}, sabiendo que pertenece a la familia de Lisp~\cite{IntroductionGuileReference}.

Con base en las propuestas anteriores se hizó un análisis comparativo (véase Anexo~\ref{Appendix1}), donde se observa que las soluciones existentes tienden a introducir nuevas fuentes de complejidad o dependen de factores externos, sin abordar las causas estructurales de la fricción en Nix; esto sugiere que un posible buen enfoque consistiría en interactuar con Nix usando un lenguaje ampliamente conocido, y TypeScript encaja muy bien en ese molde dada su popularidad~\cite{stackoverflowMostPopularTechnologies} y similitud sintáctica con Nixlang, siendo descrito en ocasiones como ``JSON con funciones'' (siendo JSON una notación usada en TypeScript)~\cite{NixdevDocumentation}.

El uso de lenguajes de propósito general para expresar dominios específicos ---conocido como lenguaje de propósito específico embebido (eDSL)~\cite{vandeursenDomainspecificLanguagesAnnotated2000}--- no es un enfoque novedoso y ha demostrado ser efectivo en contextos similares. Un caso ilustrativo es Neovim, que en 2021 introdujo un eDSL en Lua como alternativa a Vimscript~\cite{NeovimNews112021}, lo cual coincidió con un incremento notable en su interés (véase Anexo~\ref{Appendix2}). En el ecosistema de Nix, existe también una propuesta de eDSL en JavaScript; no obstante, se trata de un proyecto sin actividad reciente, limitado a un subconjunto de funcionalidades ---principalmente la creación de paquetes--- y sin evidencia empírica que respalde mejoras en la experiencia de desarrollo~\cite{burgNiJSInternalDSL2026}.

En este contexto, persiste la ausencia de una propuesta que combine un eDSL basado en un lenguaje ampliamente adoptado, con cobertura funcional más amplia de Nixlang y validación empírica de mejoras en la experiencia de desarrollo.
